% Part: intuitionistic-logic
% Chapter: tableaux
% Section: rules
%READERNOTE{SOL6-A206: उलट चिन्हांची जोडी आणि संपूर्ण शाखेची बंदता यांचे scopes वेगळे केले. नवे साक्षी-पूर्वचिन्ह किंवा त्याचा विस्तार शाखेवर आधी नसावा अशी अट स्पष्ट केली, जेणेकरून आधीच्या जगांचे अर्थनिर्धारण जपता येते.}

\documentclass[../../../include/open-logic-section]{subfiles}

\begin{document}

\olfileid{int}{tab}{rul}

\olsection{अंतःप्रज्ञावादी तर्कशास्त्रासाठीचे नियम}

$\land$ आणि $\lor$ या कारकांचे नियम नेहमीच्या विधानीय
चिन्हांकित !!{tableau}s मधील नियमांसारखेच आहेत;
फक्त त्यांना पूर्वचिन्हे जोडली जातात. प्रत्येक वेळी,
$\sFmla{S}{!A}[\sigma]$ या चिन्हांकित !!{formula}
ला नियम लावल्याने मिळणाऱ्या नव्या !!{formula}s
वरही~$\sigma$ हेच पूर्वचिन्ह असते. हे सहज लक्षात
येते: उदाहरणार्थ, $!A \land !B$ हे~$\sigma$ ने
नाव दिलेल्या जगात सत्य असेल, तर $!A$ आणि $!B$
ही दोन्ही~$\sigma$ येथे सत्य असतात (आणि त्या
जगापासून प्राप्य असलेल्या प्रत्येक जगातही सत्य
असतात). $\land$ आणि $\lor$ यांचे नियम
\olref{tab:prop-rules} मध्ये दिले आहेत.

\begin{table}
  \[\def\arraystretch{3}\begin{array}{|c|c|}
    \hline
    \AxiomC{\sFmla{\True}{!A \land !B}[\sigma]}
    \RightLabel{\TRule{\True}{\land}}
    \UnaryInfC{\sFmla{\True}{!A}[\sigma]}
    \noLine
    \UnaryInfC{\sFmla{\True}{!B}[\sigma]}
    \DisplayProof
    &
    \AxiomC{\sFmla{\False}{!A \land !B}[\sigma]}
    \RightLabel{\TRule{\False}{\land}}
    \UnaryInfC{$\sFmla{\False}{!A}[\sigma] \quad \mid \quad
      \sFmla{\False}{!B}[\sigma]$}
    \DisplayProof
    \\[2ex]
    \hline
    \AxiomC{\sFmla{\True}{!A \lor !B}[\sigma]}
    \RightLabel{\TRule{\True}{\lor}}
    \UnaryInfC{$\sFmla{\True}{!A}[\sigma] \quad \mid \quad
      \sFmla{\True}{!B}[\sigma]$}
    \DisplayProof
    &
    \AxiomC{\sFmla{\False}{!A \lor !B}[\sigma]}
    \RightLabel{\TRule{\False}{\lor}}
    \UnaryInfC{\sFmla{\False}{!A}[\sigma]}
    \noLine
    \UnaryInfC{\sFmla{\False}{!B}[\sigma]}
    \DisplayProof
    \\[2ex]
    \hline
  \end{array}\]
  \caption{$\land$ आणि $\lor$ साठीचे पूर्वचिन्हयुक्त
    !!{tableau} नियम}
  \ollabel{tab:prop-rules}
\end{table}

शाखा बंद होण्याची अट नेहमीच्या !!{tableau}s
प्रमाणेच आहे; पण येथे !!{formula}s बरोबर त्यांची
पूर्वचिन्हेही जुळली पाहिजेत. खरे तर ही अट
थोडी शिथिल करता येते: अंतःप्रज्ञावादी प्रतिमानात
!!{formula}s एकदा सत्य झाली की पुढेही सत्य राहतात.
म्हणून $!A$ हे~$\sigma$ येथे सत्य असताना
ते~$\sigma.{*}$ या प्राप्य पूर्वचिन्हावर असत्य
असू शकत नाही. म्हणून एखाद्या शाखेत पुढील
दोन्ही असतील, तर ती बंद असते:
\[
\sFmla{\True}{!A}[\sigma] \quad\text{आणि}\quad \sFmla{\False}{!A}[\sigma.{*}]
\]
येथे $\sigma$ हे एखादे पूर्वचिन्ह आणि $!A$ हे
!!{formula} आहे. उलट चिन्हे असतील, म्हणजे
शाखेत पुढील दोन्ही असतील,
\[
\sFmla{\False}{!A}[\sigma] \quad\text{आणि}\quad \sFmla{\True}{!A}[\sigma.{*}]
\]
तर सदर सूत्र-जुळणीची अट या उलट चिन्हांच्या जोडीला
तेव्हा आणि केवळ तेव्हाच लागू होते, जेव्हा $*$ ही
रिकामी क्रमिका असते. शाखा इतर कोणत्या बंदतेच्या
अटीमुळे बंद होऊ शकते, हे याने नाकारलेले नाही.

तसेच, एखाद्या शाखेत~$\sFmla{\True}{\bot}[\sigma]$
असेल, तरी ती बंद असते.

गृहीतके मांडण्याचे नियमही नेहमीच्या !!{tableau}s
प्रमाणेच आहेत; मात्र गृहीतकांसाठी आपण नेहमी
$1$ हे पूर्वचिन्ह वापरतो. (सर्व गृहीतकांसाठी
एकच पूर्वचिन्ह वापरले, तर नेमके कोणते
पूर्वचिन्ह वापरले याने फरक पडत नाही.)
म्हणून, उदाहरणार्थ,
\[
!B_1, \dots, !B_n \Proves !A
\]
असे तेव्हा आणि केवळ तेव्हाच म्हणतो जेव्हा पुढील
गृहीतकांसाठी बंद टॅब्लो असतो:
\[
\sFmla{\True}{!B_1}[1], \dots, \sFmla{\True}{!B_n}[1],
\sFmla{\False}{!A}[1].
\]

शर्त~$\lif$ साठीचे नियम अभिजात आणि मोडल
प्रकरणांपेक्षा वेगळे आहेत. $\TRule{\True}{\lif}$
हा नियम $\sFmla{\True}{!A \lif !B}[\sigma]$
असलेल्या शाखेतून दोन वेगळ्या शाखा काढतो:
एकीवर $\sFmla{\False}{!A}[\sigma.{*}]$ आणि
दुसरीवर $\sFmla{\True}{!B}[\sigma.{*}]$ जोडतो.
हा नियम लावताना $\sigma.{*}$ हे पूर्वचिन्ह
संबंधित शाखेत \emph{आधीपासूनच} आलेले
असले पाहिजे. अशा पूर्वचिन्हाला त्या शाखेवर
‘वापरलेले’ म्हणू. ($\sigma.{*}$ मध्ये
$\sigma$ स्वतःही येत असल्याने, वेगवेगळ्या
शाखांवर $\sFmla{\False}{!A}[\sigma]$ आणि
$\sFmla{\True}{!B}[\sigma]$ जोडून हा नियम
नेहमी लावता येतो.)

$\TRule{\False}{\lif}$ हा नियम
$\sFmla{\False}{!A \lif !B}[\sigma]$
असलेल्या शाखेवर
$\sFmla{\True}{!A}[\sigma.n]$ आणि
$\sFmla{\False}{!B}[\sigma.n]$ ही दोन्ही
सूत्रे जोडतो; येथे $\sigma.n$ हे त्या
शाखेसाठी नवे पूर्वचिन्ह असते: ते किंवा त्याचा कोणताही
विस्तार त्या शाखेवर आधी आलेला नसतो. पुढील नकरणाच्या
नियमांतही ‘नवे’ याच अर्थाने वापरतो.

$\lnot$ साठीचे नियम याच धर्तीवर व्याख्यायित
केले आहेत ($\lnot !A$ म्हणजे
$!A \lif \lfalse$ ही व्याख्या वापरून).

हे नियम \olref{tab:rules-lif-lnot} मध्ये दिले आहेत.

\begin{table}
  \[\def\arraystretch{3}\begin{array}{|c|c|}
    \hline
    \AxiomC{\sFmla{\True}{\lnot !A}[\sigma]}
    \RightLabel{\TRule{\True}{\lnot}}
    \UnaryInfC{$\sFmla{\False}{!A}[\sigma.{*}]$}
    \DisplayProof
    &
    \AxiomC{\sFmla{\False}{\lnot !A}[\sigma]}
    \RightLabel{\TRule{\False}{\lnot}}
    \UnaryInfC{\sFmla{\True}{!A}[\sigma.n]}
    \DisplayProof
    \\[1ex]
    \text{$\sigma.{*}$ वापरलेले आहे} & \text{$\sigma.n$ नवे आहे}\\
    \hline
    \AxiomC{\sFmla{\True}{!A \lif !B}[\sigma]}
    \RightLabel{\TRule{\True}{\lif}}
    \UnaryInfC{$\sFmla{\False}{!A}[\sigma.{*}] \quad \mid \quad
      \sFmla{\True}{!B}[\sigma.{*}]$}
    \DisplayProof
    &
    \AxiomC{\sFmla{\False}{!A \lif !B}[\sigma]}
    \RightLabel{\TRule{\False}{\lif}}
    \UnaryInfC{\sFmla{\True}{!A}[\sigma.n]}
    \noLine
    \UnaryInfC{\sFmla{\False}{!B}[\sigma.n]}
    \DisplayProof
    \\[1ex]
    \text{$\sigma.{*}$ वापरलेले आहे} & \text{$\sigma.n$ नवे आहे}\\
    \hline
  \end{array}\]
  \caption{$\lnot$ आणि $\lif$ साठीचे
    पूर्वचिन्हयुक्त !!{tableau} नियम}
  \ollabel{tab:rules-lif-lnot}
\end{table}

\end{document}
